\section{Why The Capacity Envelope Is Exhausted}
\label{sec:why-capacity-envelope-exhausted}

This section presents a closed exhaustion argument rather than a single new
theorem: it shows why every unsigned-capacity route to the terminal survival
threshold $E_b\lt T^2/(2W_-)$ fails, leaving signed residue information as
the only remaining ingredient. It holds for every nonempty conditioned chain
$5\le r_0\lt\cdots\lt r_{m-1}\lt Q$, over the same fixed eligible 2-gap-start
population and conditioned filter chain as the weighted deletion conservation
and terminal survival criterion properties above. The broader research
program contains additional mathematically proved capacity analyses. This
article uses the self-contained proofs of its load-bearing exhaustion steps in
Appendix~C; the other analyses remain in the supplementary research map.

The terminal-energy theorem of Section~5.2 reduces survival to a single inequality
on the weighted harmful-excess energy $E_b$. The most natural way to
upper-bound $E_b$ is to maximize each layer's contribution $b_i^2$ separately,
using only the arithmetic every residue histogram must satisfy. This produces
a \emph{capacity envelope}. The next six paragraphs walk through the results that
build, refine, and ultimately exhaust that envelope.

\subsection{The separate capacity envelope}
\label{sec:separate-capacity-envelope}

Each harmful residue class can hold at most $B_i$ incoming 2-gap starts, and
the two harmful classes together can hold at most $2B_i$. The sharp
capacity-envelope theorem (Appendix~C.1) proves that maximizing $b_i^2$ over
every histogram compatible with these class capacities gives the sharp
separate-layer envelope
\begin{equation*}
E_b\le\mathcal U_{\mathrm{cap}}
=\sum_i\alpha_iX_i,
\qquad
X_i=\max\bigl((\ell_i-\mu_i)^2,(u_i-\mu_i)^2\bigr),
\end{equation*}

with $\mu_i=2N_i/r_i$ and $\ell_i,u_i$ the feasible harmful-count endpoints.
The capacity-stability enlargement adds the repair
$\Gamma_{\mathrm{cap}}$, enlarging the certificate threshold to
$T^2/(2W_-)+\Gamma_{\mathrm{cap}}$. The capacity envelope is correct and
sharp \emph{for the information it retains}; the question is whether that
information is enough.

\subsection{Why capacity alone gives no positive floor}
\label{sec:capacity-no-positive-floor}

The width-floor theorem (Appendix~C.2) proves the explicit per-layer floor
\begin{equation*}
X_i\ge\frac14\min(N_i,2B_i,r_iB_i-N_i)^2.
\end{equation*}

The same property proves the obstruction: this floor vanishes at both
$N_i=0$ and $N_i=r_iB_i$. Consequently no theorem using only $r_i$ and $B_i$
can force a positive envelope, because the fully occupied profile $N_i=r_iB_i$
is positive and has zero capacity envelope. Progress requires either keeping
realized populations away from both extremes, or replacing the capacity
interval by localized residue information.

\subsection{Native-period Bessel refines but does not finish}
\label{sec:native-period-bessel}

The native-period hybrid envelope intersects Bessel's inequality over the
native prefix with the coordinate capacities via an exact greedy linear
program. This gives the hybrid envelope
$\mathcal U_{\mathrm{hyb}}\le\mathcal U_{\mathrm{cap}}$, with strict gain
exactly when a normalized prefix-capacity box exceeds an interval remainder.
The overflow quantification (Appendix~C.2) measures the gain at cut $k$ by
the normalized overflow $e_k$, which the width-floor theorem lower-bounds by
population slack. The envelope improves, but the next three subsections show
that no cut in the chain can clear the original survival threshold.

\subsection{Fixed cuts fail}
\label{sec:fixed-cuts-fail}

The fixed-seven-cut theorem (Appendix~C.3) proves that the fixed cut
immediately after filter $7$ fails: under the seven-layer density floor at
the next untouched layer (filter $11$), the hybrid envelope satisfies
\begin{equation*}
\mathcal U_2^{\mathrm{hyb}}>\frac{T^2}{2W_-}
\qquad\text{whenever }m\ge37.
\end{equation*}

The constant $37$ comes from the exact integer check
$29403\cdot275^2=2{,}223{,}601{,}875\lt37\cdot847\cdot269^2=2{,}267{,}721{,}379$.
The arbitrary-cut theorem (Appendix~C.4) generalizes this to every fixed
cut $k$:
\begin{equation*}
m>P_k(r_k-2)^2\left(1+\frac6D\right)^2
\quad\Longrightarrow\quad
\mathcal U_k^{\mathrm{hyb}}>\frac{T^2}{2W_-}.
\end{equation*}

For any fixed $k$ the right side is bounded as $Q$ grows, while $m$ tends to
infinity along any unbounded family of heads. Therefore no fixed cut can
certify survival on unbounded chains.

\subsection{Moving cuts lose their complete native blocks}
\label{sec:moving-cuts-lose-blocks}

A cut that moves outward with the head can in principle avoid the fixed-cut
obstruction. The moving-cut theorem (Appendix~C.5) proves the opposite
pressure: a threshold-clearing cut with at least one complete native block
forces
\begin{equation*}
m\lt\frac37\left(1+\frac6D\right)^2
\left(\frac{2\log H}{c}-2\right)^2,
\end{equation*}

where $H$ is the window length and $c$ is a Chebyshev-type lower-bound
constant for $\vartheta$. This bound is $O(\log^2 H)$. But the prime number
theorem gives $m=\pi(Q)-3\sim Q/\log Q$, so $m/\log^2H\sim Q/(4\log^3Q)\to\infty$.

The prime number theorem is an explicit external dependency here; the exact
logarithmic-squared inequality holds without it under five stated finite
hypotheses. The asymptotic conclusion requires PNT.

Therefore, for all sufficiently large $Q$, any threshold-clearing cut has
$M_k>H$: no complete native block remains. The incomplete-block theorem
(Appendix~C.6) then proves that the single incomplete block left at that
moving-prime scale contributes $e_k=0$---the native-capacity route gives back
exactly $\mathcal U_{\mathrm{cap}}$.

\subsection{The stability-gap repair is negligible}
\label{sec:stability-gap-negligible}

The capacity-stability enlargement had enlarged the certificate threshold by
$\Gamma_{\mathrm{cap}}$. The stability-gap theorem closes that escape: under
the seven-layer density floor at filter $7$,
\begin{equation*}
\Gamma_{\mathrm{cap}}\le\frac{25P_m}{18}\left(\frac25+\frac{3N_0}{5S}\right)^2,
\qquad
\mathcal U_{\mathrm{cap}}\ge\frac{P_mD^2}{1080}.
\end{equation*}

Prime Mertens and PNT make the stability gap eventually positive but
negligible relative to the $P_mD^2/1080$ floor. The capacity-relaxed
threshold therefore cannot rescue the separate capacity envelope on an
unbounded family.

\subsection{Verdict: signed information is required}
\label{sec:verdict-signed-information}

The capacity-plus-native-Bessel envelope cannot certify the terminal
survival threshold under the full seven-layer density floor on unbounded
chains. Every route that discards residue \emph{signs and order}---maximizing each
layer over sign-compatible histograms---has been exhausted. The shared reason
is that unsigned capacity forgets exactly the interval-order information that
controls actual harmful excess.

This sets up the next two sections. Section~7 computes the first genuine localized
saving at filter $7$ by restoring exact residue order. Section~9 then bridges that
saving to the residue-collision energy of Section~8.2 across complete old-period
copy blocks. The full self-contained
proofs of the exhaustion steps summarized in Sections~6.1--6.6 appear in
Appendix~C, so the reader can verify each constant without leaving the
article.
